\section{Requirements Analysis and Scoped Requirements Catalog for the Prototype (User Stories)}
\label{sec:requirements-analysis}

The requirements are documented in the form of user stories. According to Ott, these can serve as a catalog
of software requirements formulated in ordinary language
\cite[p.~18]{ottAppsEffektivManagen2019}. The wording follows this pattern:
\begin{quote}
  \enquote{As a <role>, I want to <goal/desire> in order to <benefit>.}
  \cite[p.~18]{ottAppsEffektivManagen2019}
\end{quote}

The requirements were extracted from the literature and expert interviews and analyzed using the approach
described in 4.1.2. The justifications refer to the sections of this thesis in which the respective
derivations have already been explained; therefore, the sources are not repeated here.

The user stories are written from the perspective of an AI facilitator who wants to proactively intervene in
the meeting using modern elements and help uncover hidden profiles. \autoref{tab:ai-facilitator-user-stories}
groups six key stories with IDs (\texttt{US1}–\texttt{US6}) so that later design and scope decisions can
clearly reference them.

\newcounter{us}
\newcommand{\usid}[1]{%
  \refstepcounter{us}%
\textbf{US\theus}\label{#1}}
\newcommand{\usref}[1]{US\ref{#1}}
\newcommand{\usblocksep}{%
  \addlinespace[0.35em]
  \midrule
\addlinespace[0.55em]}

\setcounter{us}{0}
\setlength{\tabcolsep}{6pt}
\setlength{\LTpre}{0.6em}
\setlength{\LTpost}{0.8em}
\renewcommand{\arraystretch}{1.12}
\begin{xltabular}{\textwidth}{@{}
    >{\small\raggedright\arraybackslash}p{1.15cm}
    >{\small\raggedright\arraybackslash}X
    >{\small\raggedright\arraybackslash}X
  @{}}
  \caption{Derived user stories for an AI facilitator, with references to the source chapters}
  \label{tab:ai-facilitator-user-stories}\\[0.9em]
  \toprule
  \textbf{ID} & \textbf{User Story} & \textbf{Reasoning / References} \\
  \midrule
  \endfirsthead
  \caption[]{Derived user stories for an AI facilitator (continued).}\\[0.9em]
  \toprule
  \textbf{ID} & \textbf{User Story} & \textbf{Reasoning / References} \\
  \midrule
  \endhead
  \endfoot
  \bottomrule
  \endlastfoot
  \usid{us:goal-clarity}
  & As an AI facilitator, I want to be able to clarify the meeting's agenda and objectives so that the group
  has a shared frame of reference and can reach a sound decision.
  & Structuring tasks and communicating the meeting goal are core facilitation duties and a documented
  failure mode of goal-oriented meetings
  (see \autoref{sec:facilitation}; \autoref{sec:key-challenges-in-meetings}).
  The expert interviews confirm that unclear purpose undermines facilitation quality
  (see \autoref{ch:empirical-analysis-of-facilitating-strategies}). \\
  \usblocksep
  \usid{us:quiet-voices}
  & As an AI facilitator, I want to invite people who haven’t spoken much yet to join the discussion in a
  considerate way, so that all relevant voices are heard without stigmatizing silence.
  & Unequal participation is a key meeting challenge and is linked to hierarchy and status
  (see \autoref{sec:key-challenges-in-meetings}).
  Expert facilitators invite quiet members, but warn that doing so online can stigmatize or feel like
  control
  (see \autoref{ch:empirical-analysis-of-facilitating-strategies}). \\
  \usblocksep
  \usid{us:unshared-information}
  & As an AI facilitator, I want to ask specific questions about information that hasn’t been mentioned yet
  or that doesn’t fit the picture, so that all available knowledge can be factored into the decision.
  & Groups systematically under-discuss unshared information; asking for information that does not yet
  fit the picture is a documented counter-strategy
  (see \autoref{sec:hidden-profiles}).
  The interviews describe the same practice of checking whether all relevant voices and facts are on
  the table
  (see \autoref{ch:empirical-analysis-of-facilitating-strategies}). \\
  \usblocksep
  \usid{us:hold-dissent}
  & As an AI facilitator, I don't want to resolve disagreements too quickly, so that dissent can deepen the
  exchange of information and keep better alternatives on the table.
  & Dissent improves hidden-profile solution rates; rushing through the Groan Zone closes that
  information search too early
  (see \autoref{sec:hidden-profiles}; \autoref{sec:divergent-convergent}).
  Interviewed facilitators likewise open first, then converge
  (see \autoref{ch:empirical-analysis-of-facilitating-strategies}). \\
  \usblocksep
  \usid{us:sparse-intervention}
  & As an AI facilitator, I want to intervene only sparingly and only when it clearly benefits the process,
  so as not to overshadow the human discussion.
  & A facilitator is substantively neutral and must not dominate the group
  (see \autoref{sec:facilitation}).
  Fully automatic or overly frequent intervention produces mistrust
  (see \autoref{sec:research-approaches-to-ai-supported-meeting-systems};
  \autoref{sec:interaction-paradigms}).
  Experts warn against over-facilitation and a missing sense of presence
  (see \autoref{ch:empirical-analysis-of-facilitating-strategies}). \\
  \usblocksep
  \usid{us:neutrality}
  & As an AI facilitator, I want to remain neutral regarding the content and refrain from making any
  decisions myself, so that the group retains responsibility for the outcome.
  & The facilitator is acceptable to the group, substantively neutral, and has no decision-making
  authority
  (see \autoref{sec:facilitation}).
  Interviewed facilitators invite leadership to interpret, but do not decide in the group's place
  (see \autoref{ch:empirical-analysis-of-facilitating-strategies}). \\
\end{xltabular}
\renewcommand{\arraystretch}{1}
